\section{Введение}

Оффлайн-коммуникации обеспечивают обмен сообщениями между пользователями при отсутствии доступа к сети оператора или интернету. Такая связь востребована на удалённых территориях, в экспедициях и при нарушении работы инфраструктуры. Для коротких текстовых сообщений существенны надёжность доставки и автономность узлов, тогда как высокая скорость передачи не всегда необходима. Практическую возможность маломощной системы обмена сообщениями на основе LoRa рассматривают Cardenas и соавторы~\cite{cardenas}.

LoRa задаёт физический способ передачи, но сама по себе не определяет маршрутизацию в многошаговой сети. В mesh-сети промежуточные устройства ретранслируют кадры, увеличивая достижимую область связи. Каждый дополнительный переход при этом расходует эфир и энергию. Увеличение коэффициента расширения сигнала SF изменяет чувствительность и длительность передачи, поэтому выигрыш в дальности необходимо сопоставлять с ростом занятости канала~\cite{semtech}.

Алгоритм поиска пути и протокол маршрутизации решают разные задачи. Алгоритм получает готовый граф и вычисляет путь. Протокол должен сначала обнаружить соседей, получить сведения о каналах и поддерживать их при отказах узлов. Служебные кадры используют ту же среду, что и сообщения пользователей. Поэтому малое время вычисления кратчайшего пути ещё не гарантирует высокую вероятность доставки.

Цель работы --- разработать методику оценки децентрализованной передачи сообщений, которая позволяет отдельно исследовать вычислительный уровень и сетевые затраты. Для этого используются два связанных стенда: алгоритмический и дискретно-событийный. Первый сравнивает поиск пути при известных весах, второй --- стратегии передачи при потерях и конкуренции за эфир. Связь между стендами задаётся физически интерпретируемой стоимостью ребра.

Вклад работы состоит в организации такого сравнения на согласованных входных данных и в явном разграничении оптимального одноадресного маршрута и стратегии с несколькими потенциальными ретрансляторами. Настоящая версия статьи содержит постановку задачи и описание методов. Численные результаты и итоговые рекомендации по выбору протокола в неё не включены.
